\section{Resultados}

Na figura 1, encontra-se a montagem experimental realizada, de onde se extraíram-se os valores da tensão nos terminais do condensador. É constituída por uma fonte de tensão de $\qty{15}{\volt}$, duas resistências $R_1 = 1,5$ k$\Omega$ e $R_2 = 1$ k$\Omega$ associadas em série, um condensador $ C = \qty{10}{\milli\farad}$, associado em paralelo com as resistências e dois interruptores, $S_1$ e $S_2$.

\begin{figure}[!ht]
	\centering
	\begin{circuitikz}[american voltages]
		\draw (0,0) to [battery1, v<={$V=\qty{15}{V}$}, invert] ++(0,2)
		to [nos=$S_1$, n=S1] ++(2,0)
		node[ocirc] at (S1.e) {}
		node[ocirc] at (S1.w) {}
		to [R={$R_1=\qty{1.5}{k\ohm}$}] ++(2,0)
		node (top) {}
		to [nos={$S_2$}, n=S2] ++(2,0)
		node[ocirc] at (S2.e) {}
		node[ocirc] at (S2.w) {}
		node (V+) {}
		to [C={$C=\qty{10}{\milli\farad}$}, n=C1] ++(0,-2)
		node (V-) {}
		-- (0,0);

		\draw (top) to [R, l_={$R_2=\qty{1}{k\ohm}$}] ++(0,-2);

		\draw (V+) -- ++(3,0)
		to [voltmeter] ++(0,-2)
		-- (V-);

	\end{circuitikz}
	\caption{Montagem experimental}
	\label{circ:montagem}
\end{figure}

\subsection{Carga do Condensador}
Foi realizada a montagem do circuito ilustrado na Figura 1, recorrendo a 2 resistências $R_1=1.5ohm$
\begin{table}[!ht]
	\centering
	\begin{tabular}{|c|c|c|c|c|c|c|c|}
		\hline
		$t \pm \Delta t \quad (\unit{s})$ & 6    & 12   & 18   & 24   & 30   & 36   & 42   \\
		\hline
		$V \pm \Delta V \quad (\unit{V})$ & 2,35 & 4,14 & 5,15 & 5,55 & 5,72 & 5,81 & 5,85 \\
		\hline
	\end{tabular}
	\caption{Carga}
	\label{tab:carga}
\end{table}

\subsection{Descarga do Condensador}
-----

\begin{table}[!ht]
	\centering
	\begin{tabular}{|c|c|c|c|c|c|c|c|}
		\hline
		$t \pm \Delta t \quad (\unit{s})$ & 6    & 12   & 18   & 24    & 30    & 36    & 42    \\
		\hline
		$V \pm \Delta V \quad (\unit{V})$ & 3,89 & 2,40 & 1,35 & 0,820 & 0,533 & 0,327 & 0,202 \\
		\hline
	\end{tabular}
	\caption{Descarga}
	\label{tab:descarga}
\end{table}


